%% kapa: exam=1 round=30 year=2019 period=1 form=A subject=economics subject_ko=경제학원론 seq=32 number=72
%% source: Q-Net 감정평가사 기출문제 (cst00309) artlSeq=5205205 file=extracted/1008/5205205/2019년 제30회 감정평가사 1차시험 문제지/2019년 제30회 감정평가사 1차시험 문제지/2019년도 제30회 감정평가사 1차 1교시 A형.hwp
%% conversion: kordoc 4.24.51 -> tools/gen1.py (auto); review=proofread-sonnet; flags=none
\begin{question}{72}
다음 거시경제모형에서 생산물시장과 화폐시장이 동시에 균형을 이루는 소득과 이자율은? (단, ${ C }$는 소비, ${ Y }$는 국민소득, ${ I }$는 투자, ${ G }$는 정부지출, ${ T }$는 조세, ${ r }$은 이자율, ${ MD }$는 화폐수요, ${ MS }$는 화폐공급이다. 물가는 고정되어 있고, 해외부문은 고려하지 않는다.)
\begin{framedtext}
○ ${ C=20+0.8(Y-T)-0.5r }$ \par
○ ${ I=50-9.5r }$ \par
○ ${ G=50 }$ \par
○ ${ T=50 }$ \par
○ ${ MD=50+Y-50r }$ \par
○ ${ MS=250 }$
\end{framedtext}
\begin{choices}
  \item 200, 1
  \item 200, 2
  \item 250, 1
  \item 300, 1
  \item 300, 2
\end{choices}
\end{question}
\answer{5}
